%==============================================================================
% Cover letter — IJHMT submission
% PINN-ONB01 manuscript
%==============================================================================
\documentclass[11pt,a4paper]{letter}

\usepackage[T1]{fontenc}
\usepackage{lmodern}
\usepackage[utf8]{inputenc}
\usepackage[margin=2.5cm]{geometry}
\usepackage{microtype}
\usepackage{hyperref}
\hypersetup{colorlinks=true, linkcolor=blue, urlcolor=blue}

\signature{Jaeseon Lee\\
Corresponding author\\
Innovative Thermal Engineering Laboratory\\
Department of Mechanical Engineering\\
Ulsan National Institute of Science and Technology\\
50 UNIST Rd., Ulsan 44919, Republic of Korea\\
\href{mailto:JaeseonLee@unist.ac.kr}{JaeseonLee@unist.ac.kr}\\
Tel.\ +82-52-217-2342}

\address{Jaeseon Lee\\
Innovative Thermal Engineering Laboratory\\
Department of Mechanical Engineering\\
Ulsan National Institute of Science and Technology\\
50 UNIST Rd., Ulsan 44919, Republic of Korea}

\longindentation=0pt

\begin{document}

\begin{letter}{The Editor-in-Chief\\
\emph{International Journal of Heat and Mass Transfer}\\
Elsevier}

\opening{Dear Editor,}

We are pleased to submit our manuscript entitled \emph{``Surface-Conditioned
Physics-Informed Neural Network for Pool-Boiling Onset of Nucleate
Boiling''} for consideration for publication in the \emph{International
Journal of Heat and Mass Transfer}.

The accurate prediction of the onset of nucleate boiling (ONB) on
engineered surfaces remains an open problem. Existing analytical
correlations rely on idealised cavity-distribution assumptions and treat
surface character through a single scalar descriptor. Their residual
scatter is 20--30\%, and they cannot accommodate roughness control,
wettability tailoring, biphilic patterning, or nanostructured coatings
in a unified framework.

This work makes the following contributions:
\begin{enumerate}
  \item A curated and standardized pool-boiling dataset of 1{,}361
        boiling-curve points and 82 ONB labels across four fluids (water,
        R-134a, R-123, FC-77), drawn from seven primary literature sources
        and accompanied by 49 surface descriptor cards.
  \item A surface-conditioned PINN architecture (24{,}005 parameters) in
        which a compact encoder modulates the temperature-field network
        through feature-wise linear modulation (FiLM). A single backbone
        represents qualitatively distinct surface families.
  \item A five-term composite loss that combines the conduction PDE
        residual, the natural-convection boundary condition, experimental
        ONB data, a Hsu-type soft nucleation constraint, and
        monotonicity regularizers.
  \item A Hsu-based inverse-problem study that recovers active cavity
        radii and uncovers a Simpson-type reversal in the
        roughness--cavity correlation across fluids.
  \item Code verification (four-of-four Level-1 tests, maximum relative
        error 0.026\%) and a Level-3 physical-consistency battery (eight
        of nine tests passing). The model attains an ONB RMSE of
        3.42~K with $R^{2}=+0.44$ on $n=77$ held-out points, a 53\%
        reduction over the best classical correlation and 65--67\% on
        the refrigerant subset.
\end{enumerate}

To our knowledge this is the first study to combine pool-boiling ONB,
explicit surface-modification conditioning, multi-source cross-fluid
generalization, and inverse cavity-radius recovery within a single
physics-informed framework. We believe these advances are directly
relevant to the readership of the journal and to the design of
engineered boiling surfaces for nuclear safety systems, electronics
thermal management, and refrigeration evaporators.

This manuscript has not been published previously, is not under
consideration for publication elsewhere, and its submission has been
approved by all authors. We declare no conflicts of interest. The
curated dataset, model code, and analysis scripts will be released as
open-source supplementary material upon acceptance.

We thank you for your time in considering our submission and look
forward to your reviewers' comments.

\closing{Yours sincerely,}

\end{letter}
\end{document}
